\section{Where the Error Lives}
\label{sec:where}

\subsubsection{Field altitude and variable identity.}
Near-surface means surface fields plus 1000/850\,hPa (13 channels); upper-air
means 500/50\,hPa (6). These refer to field altitude, not ground elevation;
\texttt{tcwv} spans the atmospheric column and belongs to neither group.
The reductions are $G=32.2\%$ versus $3.5\%$ and $\bar G=27.5\%$ versus
$1.5\%$. Their ratios, 9.1 and 18.6, compare gains rather than compression.
Surface-pressure nRMSE falls from 0.0990 to 0.0249; the 500\,hPa wind errors
are unchanged to four decimals. The largest upper-air gain is 7.4\% for relative
humidity at 500\,hPa.

Across 1000, 850 and 500\,hPa, $u$ gains are 40.5\%, 20.1\% and approximately
zero; $v$ gives 40.7\%, 19.6\% and approximately zero. From 850 to 500\,hPa,
$t$ falls from 12.3\% to 0.1\%, $r$ from 15.1\% to 7.4\%: variable identity alone
does not explain the split.
\begin{table}[t]
\centering
\caption{Full-2018 frozen-checkpoint evaluation (1{,}460 frames). Before/after
nRMSE averages 13 near-surface or 6 upper-air direct channels; both exclude
column-integrated \texttt{tcwv}. $G$ is their relative reduction. CR uses measured entropy-coded bytes: R is CRA5 reference,
M is matched mean rate, and N is non-matched. e995/e99/e97 denote MPCA+PCA.
Each system reuses its own previously trained weights; MPCA-only was trained on
stage-1-only reconstruction and is not retrained for this table. Classical rows
use actual quantised reconstruction at the original fixed steps.}
\label{tab:asymmetry}
\begin{tabular}{lllrr@{\hspace{2em}}r}
\toprule
System & CR/status & Group & Before & After & $G$ (\%) \\
\midrule
CRA5 & 471.6 R & Near & 0.09820 & 0.06660 & 32.18 \\
 & & Upper & 0.06151 & 0.05934 & 3.52 \\
MPCA-only & 473.1 M & Near & 0.12913 & 0.08207 & 36.44 \\
 & & Upper & 0.12289 & 0.09690 & 21.14 \\
e995 & 1333 N & Near & 0.17983 & 0.14221 & 20.92 \\
 & & Upper & 0.15849 & 0.13392 & 15.50 \\
e99 & 2000 N & Near & 0.19156 & 0.15536 & 18.89 \\
 & & Upper & 0.17149 & 0.14686 & 14.36 \\
e97 & 5093 N & Near & 0.22290 & 0.19081 & 14.39 \\
 & & Upper & 0.20767 & 0.18505 & 10.89 \\
\bottomrule
\end{tabular}
\end{table}

\subsubsection{Classical controls and architecture.}
Near/upper $G$ ratios are 9.14 for CRA5, 1.72 for matched MPCA-only, and
1.32--1.35 for non-matched MPCA+PCA (Table~\ref{tab:asymmetry}). Distinct
checkpoint weights share an architecture and residual objective. This supports
a compressor-dependent pattern for this corrector, but its local receptive field
may preferentially repair surface structure while missing broader upper-air
errors. The MPCA architecture sweep does not isolate CRA5 receptive-field effects;
we therefore test a fixed parameter-matched intervention below.

\subsubsection{Parameter-matched receptive-field ablation.}
DilatedUNet2 changes only the two bottleneck convolutions to dilation/padding
2 and 3, preserving all other operations, shapes and 1{,}646{,}548 parameters.
The same CRA5 residuals and training recipe select epoch 125 on 2017, stopping
at 137 before testing. This fixed ablation is not architecture search.
Training/validation takes 2{,}278\,s. Both models use 19.86\,G counted
FLOPs/frame; same-run batch-8 inference averages 1.316/1.315\,ms/frame.

Including pooling, bilinear resizing and skips, central theoretical support
grows $72\times72\to96\times96$ cells (unclipped maximum spans 76\,$\to$\,100
per axis). Gradient diagnostics of $f(x)$ use eight fixed 2017 inputs, all 20
outputs and three positions: 480 probes/model. Squared gradients are summed
over input channels and each probe normalised. Central mean 90\%-energy Chebyshev radius decreases
3.41\,$\to$\,2.78 cells. Across all probes, mean Dilated energy outside the
original support is only 0.0033\%. Expanded theoretical support does not imply broad effective use.

Overall correction is slightly worse (Table~\ref{tab:outtime}). In 2018/2022,
upper $G$ rises 0.30/0.31 percentage points (paired 95\% intervals [0.27,0.33]/[0.28,0.34]),
while near $G$ falls 0.57/0.58 points. Near/upper ratios remain 8.28/8.24;
differences are 27.79/28.30 points. $z500$ worsens ($G=-4.46\%/-5.12\%$),
and upper $\bar G$ falls to 0.94\%/1.01\%. This ablation retains the split,
but concentrated gradients and one seed leave broader architecture effects open.

\subsubsection{Low-capacity affine control.}
One $3\times3$ convolution (20\,$\to$\,20, bias, no activation) has 3{,}620
parameters: \emph{not parameter-matched}. The same recipe selects epoch 125 on
2017, stopping at 137; frozen direct19 gains are only 1.29\%/1.31\% in
2018/2022. Its 2018 upper $\bar G=-0.91\%$. Different capacity and context
prevent attributing the U-Net gap solely to nonlinearity.

\subsubsection{Spatial diagnostic.}
For a $5\times5$ box smoother $S_5$, fine-scale error fraction
$\mathrm{Var}(e-S_5(e))/\mathrm{Var}(e)$ is 57.6\%/27.6\% near/upper for CRA5,
versus 22.2\%/5.5\% for unquantised e995. This diagnostic does not establish a mechanism.
